\documentclass[10pt,a4paper]{amsart}
\usepackage[margin=19mm]{geometry}
\usepackage{amsmath,amssymb,mathtools}
\usepackage[scr=rsfs,cal=euler]{mathalfa}
\usepackage{xfrac}
\usepackage[unicode,hidelinks,bookmarks=false]{hyperref}
\newcommand{\ie}{i.e.\ }
\newcommand{\cf}{cf.\ }
\newcommand{\resp}{respectively\ }
\newcommand{\Subsection}{\subsection}
\providecommand{\nobreakdash}{\nobreak\hbox{-}}
\setlength{\emergencystretch}{2em}
\multlinegap0pt
\usepackage{bm}
\usepackage{bbm}
\usepackage{dsfont}
\usepackage{xspace}
\usepackage{cancel}
\usepackage{mathtools}
\usepackage{amsthm}
\makeatletter
\@ifundefined{theorem}{\newtheorem{theorem}{Theorem}}{}
\makeatother
\makeatletter
\@ifundefined{poc}{\newenvironment{poc}{\begin{proof}[Proof of
    Claim]\renewcommand*{\qedsymbol}{$\blacksquare$}}{\end{proof}}}{}
\makeatother
\title[Dirac's theorem and the switch geometry of perfect matchings]{Dirac's theorem and the switch geometry of perfect matchings}
\author{Ross J. Kang, Cl\'ement Legrand-Duchesne}
\date{Source: arXiv:2604.17911v1 (CC BY 4.0)}
\begin{document}
\maketitle

\noindent\emph{Source and licence.} Dirac's theorem and the switch geometry of perfect matchings. Version arXiv:2604.17911v1 (CC BY 4.0), distributed under the Creative Commons Attribution 4.0 International licence, as recorded in the arXiv licence metadata for this exact version and shown on the abstract page https://arxiv.org/abs/2604.17911v1. Full rights notice: licenses/arxiv-2604.17911-CC-BY-4.0.txt.\par\smallskip

\noindent\textit{Editorial scope.} Counted unit: the Theorem whose source label is \texttt{thm:isolated\_k=2} in the article (source0.tex lines 2741--2750, source label \texttt{thm:isolated\_k=2}), delivered together with the article's own complete original proof of it (source0.tex lines 2751--2797, one \texttt{proof} environment of 47 lines). No mathematics is written, repaired or re-derived: statement and proof are the article's own text line for line. Required context retained uncounted: thm:isolated\_Caccetta (lines 2672--2677). Every definition, display and label of those lines is retained unchanged. Omitted from this package and not redistributed: 1--2671, 2678--2740, 2798--2893. Nothing inside a retained excerpt is omitted or altered: every retained body is delivered exactly as the article's lines are written, so no omission receipt of any kind accompanies this package. Citations: the retained excerpts carry no citation command at all, so no bibliography entry of the article is transcribed and no reference is redistributed. Reference adaptation: none was needed and none was made; every reference inside the delivered unit resolves inside it, so no unresolved reference or citation is produced. Printed slips kept literal: no printed word, symbol, hypothesis or number of the counted statement or of its proof was altered. Layout and class: the article's own class is \texttt{article} with packages including graphicx, amsfonts, amsmath, amssymb, amsthm, mathtools, fontenc, babel, inputenc, geometry, url, hyperref, thm-restate, tikz; the wrapper is the lane's pinned \texttt{amsart} editorial wrapper at 10pt on A4, which restates the theorem-like environments the retained text needs and adds the article's own macro definitions that the retained text needs (none), with extra wrapper packages (bm, bbm, dsfont, xspace, cancel, mathtools). The delivered numbering is the unit's own. Assets: no retained span contains a figure environment, a graphics-inclusion command, a PDF-page inclusion, a table or a TikZ picture; no image file is copied into this package, the package declares no asset dependency and no image file is redistributed with it. Licence: version arXiv:2604.17911v1 of this article is distributed under the Creative Commons Attribution 4.0 International licence, as recorded in the arXiv licence metadata for this exact version and shown on the abstract page https://arxiv.org/abs/2604.17911v1. A full rights notice is delivered at licenses/arxiv-2604.17911-CC-BY-4.0.txt. Characters: every retained excerpt is pure ASCII, so the delivered engine, XeLaTeX with its default Latin Modern text font, reports no missing character.
\begin{theorem}\label{thm:isolated_Caccetta}
  Let $k \ge 2$ and $n$ be an even integer. Let $d_k(n)$ be the largest integer such
  that there exists an $n$-vertex graph with minimum \emph{outdegree} $d_k(n)$ and no
  directed cycle of length at most $k$. Then
  $$\delta^{\mathbf{no iso}}_k(n) \le d_k(n)+2.$$
\end{theorem}

\begin{theorem}\label{thm:isolated_k=2}
  Let $G$ be $n$-vertex graph (or respectively a balanced bipartite graph on
  $2n$ vertices). If $\delta(G) \ge n/2+1$ (respectively
  $\delta(G) \ge \lfloor (n+3)/2\rfloor$), then for every perfect matching $M$
  of $G$ and every $e \in E(M)$, there exists a $2$-switch that contains $e$.\\
  On the other hand, $\mathcal H_2(G)$ may have isolated vertices if
  $\delta(G) \le n/2 -2$ (respectively $ \lfloor (n+1)/2\rfloor$). In other words,
  $$ \frac{n}2-1 \le \delta^{\mathbf{no iso}}_2(n) \le \frac{n}2+1  \quad \text{ and } \quad
  \delta^{\mathbf{no iso}}_{2,\mathrm{bip}}(n) = \left\lfloor \frac{n+3}2 \right\rfloor.$$
\end{theorem}

\begin{proof}
  We first prove the existence of such a switch if the minimum degree is sufficiently high. Let $M$ be a perfect matching of $G$ and $uv$ be an edge of $M$. Let $V' = V(G) \setminus \{u,v\}$, $A$ be the subset of vertices of
  $V'$ that are matched in $M$ to a neighbour of $u$, and $B = N(v) \cap V'$
  \footnote{If $G$ is a balanced bipartite graph, then let
    $V' = U \setminus \{u\}$, where $U$ is the half of the bipartition
    containing $u$, and let $A$ and $B$ be defined identically. Then $A$ and $B$ have
    size at least $\lfloor (n+1)/2\rfloor$, so $|A| + |B| \ge n > |V'|$ and by the
    pigeonhole principle, $A$ and $B$ intersect.}. We have
  $|A| \ge \deg(u) -1 \ge n/2$ and $|B| \ge n/2$. By the pigeonhole principle, $A$
  and $B$ intersect, so there exists a $2$-switch that contains $uv$.

  Conversely, we construct graphs with high minimum degree and an isolated
  matching. To do so, we first construct an auxiliary oriented graph $H_p$ on
  $p$ vertices with maximal semidegree
  $\delta(H_p) := \max_{u_i} \max \{\delta^+(u_i), \delta^-(u_i)\}$. The
  semidegree of an oriented graph on $p$ vertices is at most
  $\lfloor (p-1)/2 \rfloor$. This bound is attained by taking $H_p$ to be the
  oriented graph such that there is an arc from $u_i$ to $u_j$ is
  $j-i \in \{1, \dots, \lfloor (p-1)/2 \rfloor\} \pmod p$. By
  Theorem~\ref{thm:isolated_Caccetta}, there exists a balanced bipartite graph
  $G_n^{(\mathcal B)}$ on $2n$ vertices with minimum degree
  $\lfloor (n-1)/2 \rfloor +1 = \lfloor (n+1)/2\rfloor$, such that
  $\mathcal H_2(G_n^{(\mathcal B)})$ contains an isolated vertex.

  We now give the construction for general graphs. Let $n$ be an even integer and $p = n/2$
  and consider the graph $G_n$ on the vertex set
  $\{a_1, \dots a_{p}\} \cup \{b_1, \dots b_p\}$ built from $H_p$ as
  follows. For each $i$, connect $a_i$ with $b_i$. For each vertex
  $u_i \in V(H_p)$, we partition $N^+(u_i)$ into $A_i \cup B_i$ such that
  $||A_i|- |B_i||\le 1$. We connect $a_i$ to $a_j$ and $b_j$ for all
  $u_j \in A_i$, and $b_i$ to $a_j$ and $b_j$ for all $u_j \in B_i$. The graph
  $G_n$ constructed this way has minimum degree equal to
  $1 + 2\lfloor \delta^+(H_p)/2\rfloor + \delta^-(H_p)$. Moreover, the matching
  $M = \{a_ib_i \colon i \in [n/2]\}$ is isolated in $\mathcal H_2(G_n)$: since $H_p$ is
  an oriented graph, for each distinct $i$ and $j$, the subgraph of $G_n$
  induced by $\{a_i,b_i,a_j,b_j\}$ is isomorphic to a triangle with an edge
  attached to one of the vertices, so the edges $a_ib_i$ and $a_jb_j$ cannot be
  part of a $2$-switch. The minimum degree of $G_n$ is
  \begin{align*}
    \delta(G_n) &= 1 + 2 \left\lfloor \frac{p-1}4 \right \rfloor +  \left\lfloor
                  \frac{p-1}2 \right \rfloor
    = 2\left\lceil \frac{p}4 \right \rceil +  \left\lceil
      \frac{p}2 \right \rceil -2\\
    &= 2\left\lceil \frac{n}8 \right \rceil +  \left\lceil
      \frac{n}4 \right \rceil -2 \ge n/2 -2. \qedhere
  \end{align*}
\end{proof}

\end{document}
